We used two-component Normal--Student-$t$ and Normal--skew-normal mixtures in
dimensions $p=2$ and $p=5$, with Euclidean location separations representing
separated, moderate-overlap, and strong-overlap regimes.  Balanced designs used
$n=500$ and $n=2000$; additional experiments considered weights $(0.2,0.8)$,
three heterogeneous components, and $p=10$.  There were 1000 replications per
primary cell and 500 in the dimension stress cells, for 31,000 attempted fits.
Bandwidth multipliers $0.75$, $1$, and $1.5$ were evaluated without refitting.
All random seeds, including failed replications, were retained.

Figure~\ref{fig:simulation-profiles} shows the principal two-dimensional
designs before aggregation over regimes.  The correction has little effect in
the separated cells, but the gap between the methods grows as the active
boundary moves into a region of greater probability mass.

\begin{figure*}[t]
\centering
\includegraphics[width=\textwidth]{figures/simulation_profiles.pdf}
\caption{Standard-error calibration and 95\% coverage in the principal
two-dimensional simulations. Shapes distinguish separation regimes; solid
blue and dashed grey curves denote stabilized boundary-aware and naive
inference. Each point is a coordinate median over usable replications.}
\label{fig:simulation-profiles}
\end{figure*}

The estimand was the population classification maximum-likelihood functional,
not the generating ordinary-mixture parameter.  For each of 18 scenarios we
used 10 independent reference samples of size 100,000, reduced to five for the
initial $p=10$ stress calculation.  All reference fits succeeded.  Across the
18 targets the largest coordinate Monte Carlo standard error was 0.033 and the
largest range across independent reference fits was 0.322.  No selected target
fit violated the prespecified admissibility diagnostics.

Table~\ref{tab:simulation-main} reports coordinate-median coverage and the
coordinate-median ratio of estimated standard error to empirical standard
deviation.  The stabilized method uses the central bandwidth multiplier.  The
correction is most consequential under overlap and imbalance.  For example, in
the Normal--Student-$t$, $p=2$, moderate-overlap design with $n=500$, median coverage
increased from 0.886 to 0.925 and the median standard-error ratio from 0.758 to
0.920.  Under imbalance the corresponding coverage increased from 0.834 to
0.895.  Under separation and $n=2000$, both methods are closer, as expected.

\begin{table}[t]
\centering
\caption{Selected simulation results. Entries are coordinate-median estimated
standard error divided by empirical standard deviation, followed in
parentheses by coordinate-median 95\% coverage.}
\label{tab:simulation-main}
\begin{tabular}{llrrcc}
\toprule
Families & Regime & $p$ & $n$ & Naive & Stabilized \\
\midrule
Normal--Student-$t$ & moderate & 2 & 500  & 0.758 (0.886) & 0.920 (0.925) \\
Normal--Student-$t$ & moderate & 2 & 2000 & 0.827 (0.896) & 0.994 (0.940) \\
Normal--Student-$t$ & imbalanced & 2 & 500 & 0.615 (0.834) & 0.844 (0.895) \\
Normal--Student-$t$ & strong & 2 & 500 & 0.491 (0.750) & 0.675 (0.844) \\
Normal--Student-$t$ & separated & 5 & 2000 & 0.956 (0.939) & 1.018 (0.950) \\
Normal--skew-normal & strong & 2 & 2000 & 0.573 (0.868) & 0.762 (0.922) \\
Normal--skew-normal & moderate & 5 & 2000 & 0.892 (0.933) & 0.994 (0.943) \\
\bottomrule
\end{tabular}
\end{table}

Table~\ref{tab:simulation-quantiles} aggregates all retained coordinate-level
primary results rather than selecting favourable coordinates.  Coordinate-level
results, including the difficult coordinates, are reported in the supplement;
no difficult coordinate was removed from the analysis.

\begin{table}[t]
\centering
\caption{Across-coordinate distribution of calibration results in the primary
study.  Entries are the 25th percentile, median, and 75th percentile across all
scenario--sample-size--coordinate cells; corrected results use the central
bandwidth.}
\label{tab:simulation-quantiles}
\begin{tabular}{lrrrrrr}
\toprule
& \multicolumn{3}{c}{SE/empirical SD} & \multicolumn{3}{c}{Empirical coverage}\\
Method & Q25 & Median & Q75 & Q25 & Median & Q75\\
\midrule
Naive & 0.654 & 0.867 & 0.952 & 0.878 & 0.917 & 0.935\\
Corrected, stabilized & 0.780 & 0.991 & 1.032 & 0.911 & 0.937 & 0.947\\
\bottomrule
\end{tabular}
\end{table}

Figure~\ref{fig:simulation-calibration} summarizes every primary scenario.
Nearly all points lie above the equality line: the correction increases both
the estimated-standard-error ratio and coverage, generally toward the ideal
horizontal reference rather than indiscriminately beyond it.

\begin{figure*}[t]
\centering
\includegraphics[width=\textwidth]{figures/simulation_calibration.pdf}
\caption{Scenario-level comparison of naive and stabilized boundary-aware
inference across the complete primary study. Point colour denotes the family
design and point size distinguishes larger sample sizes. The dashed diagonal
denotes equality; dotted lines mark ideal SE calibration and nominal coverage.}
\label{fig:simulation-calibration}
\end{figure*}

Across the primary study, 4214 scenario--replication cases had at least one raw
nonpositive curvature result among the three prespecified bandwidths.  This is
why stabilization is reported as a diagnostic rather than hidden as a purely
numerical adjustment.  Raw correction frequently produced a nonpositive estimated curvature, with the
frequency increasing under strong overlap and in higher dimension.  In
contrast, the stabilized correction was computable whenever the underlying
fit and naive information calculation succeeded.  Stabilization was active in
10.5\% of Normal--Student-$t$, $p=2$, moderate-overlap replications at $n=500$, falling
to 3.4\% at $n=2000$; in the separated $p=2$ cells it fell from 5.1\% to 1.0\%.
Under the regular theory, stabilization frequency should decrease with sample
size; these paired simulation cells show precisely that tendency.

Normal--Student-$t$ fitting succeeded in every primary replication.  Operational fit
failure for models involving a skew-normal component ranged from zero to 7.6\%,
with the maximum occurring in the $p=10$ stress design.  We did not replace or
discard these seeds.  Conditional summaries use all replications for which the
base fit and naive information were available, and failure rates are reported
separately.

A post-confirmatory convergence extension used 500 replications at $n=8000$
for the two $p=2$ strong-overlap models and at $n=4000$ for the two
three-component models.  Stabilized median coverage reached 0.940 for the
Normal--skew-normal strong-overlap case and 0.942 for the $p=5$ three-component
case.  The Normal--Student-$t$ strong-overlap case improved to 0.923, compared with
0.721 for naive inference.  The $p=2$ three-component case remained difficult,
with stabilized median coverage 0.863.  This last result illustrates slow
finite-sample convergence when several decision faces and family-assignment
uncertainty interact; it is reported as a limitation rather than removed from
the design.

Figure~\ref{fig:extension-boundary} places this extension beside a separate
boundary-estimator convergence experiment. The latter shows decreasing RMSE
from sample size 250 to 16,000, directly checking the feasible surface-moment
estimator independently of the Wald calculations.
The supplement additionally uses all replication-level estimator vectors to
check standardized root-$n$ quantiles, Kolmogorov distance from Gaussianity,
scaled RMSE, and scaled bias. Thus consistency, rate, asymptotic normality,
variance calibration, coverage, feasible-curvature convergence, stabilization,
and bootstrap behaviour each receive a distinct empirical diagnostic.

\begin{figure*}[t]
\centering
\includegraphics[width=\textwidth]{figures/extension_boundary_convergence.pdf}
\caption{Large-sample and estimator-convergence diagnostics. Left:
coordinate-median coverage in four post-confirmatory cells; segments join
naive and corrected results. Right: RMSE of a scalar boundary-curvature
functional, with an $n^{-1/2}$ reference slope.}
\label{fig:extension-boundary}
\end{figure*}
